 \documentclass[11pt]{article}
\usepackage{amsmath,amsfonts}
\usepackage{graphicx}
%\usepackage{xcolor}
%\definecolor{gray}{rgb}{.75,.75,.75}
%\usepackage[landscape]{geometry}
%\usepackage{hyperref}
%\usepackage[bookmarks=true, pdfborder={0 0 .5}, urlbordercolor=gray]{hyperref}

\renewcommand{\thefootnote}{\fnsymbol{footnote}}


\addtolength{\topmargin}{-.6in}
\addtolength{\textheight}{1.0in}
\addtolength{\oddsidemargin}{-.35in}
\addtolength{\textwidth}{.7in}
\pagestyle{empty}

\newcommand{\ds}{\displaystyle}
\newcommand{\ts}{\textstyle}

\newcommand{\Z}{\mathbb{Z}}
\newcommand{\R}{\mathbb{R}}
\newcommand{\C}{\mathbb{C}}
\newcommand{\EE}{\mathcal{E}}
\newcommand{\tZ}{\tilde{\Z}}

\newcommand{\Sym}{\mathrm{Sym}}

\renewcommand{\circle}[1]{{\bigcirc\hspace{-.75em}#1}}

\begin{document}

\footnote[0]{Notes by Peter Magyar \ {\tt magyar@math.msu.edu}}
\vspace{-1em}

\centerline{\bf Math 133\hspace{1.5em} \hfill{\Large\bf
Parametric Curves}\hfill 
Stewart \S10.1}
\vspace{1em}

\noindent 
{\bf Back to pictures!} Recall our four conceptual levels, 
points of view on mathematics: physical, geometric, numerical, algebraic.
The physical view is Applied Mathematics, including
engineering and the hard sciences: distilling key features of the real world
into highly simplified mathematical models allows us
to predict and control what happens (if nothing goes wrong).
The numerical view, called Analysis,
deals with approximations, error-control, and convergence of limits.
It guides our reasoning through the perplexities of 
infinite processes, but it is a special taste even among mathematicians.  
Algebra, my favorite, is concerned with formulas
to consisely construct and transform complicated quantities by 
symbolic operations, often giving amazingly simple answers.

But deep down, what we really love in math is Geometry: pictures! 
In this section, we will learn to handle the simplest geometric objects: curves.
So far, we have dealt with curves as graphs of functions $y=f(x)$,
in which we imagine the independent variable $x$ moving horizontally along its axis while $f(x)$ controls the height $y$.
%This has two drawbacks: it is difficult to describe a curve
%which does not satisfy the vertical line test, 
%so is not the graph of a single function;
%and it does not describe the motion of a particle 
%along a curve.
\\[1em]
{\bf Parametric lines.}
A more general model for a curve is to consider it as the path of a particle moving in the plane in any fashion. We specify its coordinates as functions of time:
at time $t$, the particle is at position $(x(t),y(t))$. 
We call the controlling variable $t$ the {\it parameter}, 
and the trajectory traced out is a {\it parametric curve}.  

Any graph $y=f(x)$ can immediately be written parametrically as
 $(x(t),y(t))=(t, f(t))$, meaning at time $t$ the particle is above the point 
$x=t$ with height $y=f(t)$.
\\[.8em]
{\sc example:} Suppose a particle starts at time $t=0$ 
at the point $(x(0),y(0)) = (1,2)$ , 
and moves with constant velocity until time $t=1$ to the point
$(x(1),y(1)) = (4,6)$. The horizontal velocity is $\frac{\Delta x}{\Delta t}=\frac{4-1}{1-0} = 3$, 
the vertical velocty is $\frac{\Delta x}{\Delta t}=\frac{6-2}{1-0}=4$,\footnote{
The two-dimensional velocity $\vec v = (3,4)$ is called the {\it velocity vector},
pictured as an arrow from $(0,0)$ to $(3,4)$, or any arrow with the same direction
and length, such as from the start point $(1,2)$ to the end point $(4,6)$.
It points in the direction of motion,
and its length is the overall speed $|\vec v|=\sqrt{3^2+4^2} = 5$.
} 
and the position at time $t$, moving along the thick line segment below. will be:
$$
(x(t), y(t)) \ = \ (1,2) + (3t,4t) \ = \ (1+3t,\, 2+4t)\quad\text{for}\quad t\in[0,1]\,,
$$
\\[-.7em]
\centerline{
\includegraphics[width=1.5in]
{10-1-Line-Segment.png} 
}
\vspace{.0em}

\noindent
If we keep the same velocity for all $t\in(-\infty,\infty)$, we get the 
thin infinite line. 

\newpage

\noindent
Given any parametric curve, writing it in terms of an equation in $x$ and $y$
(with no $t$) is called {\it deparametrizing}:
in this case, we want the graph of a linear function $y=mx+b$.
A general method is to solve for $t$ in terms of $x$,
then plug in to the equation for $y$:
$$
\left\{ \begin{array}{c}
x=1+3t\\
y=2+4t
\end{array}\right.
\quad\Longrightarrow\quad
\left\{ \begin{array}{c}
t=\tfrac13(x{-}1)\\[.4em]
y=2+4\!\left(\tfrac13(x{-}1)\right)
\end{array}\right.
\quad\Longrightarrow\quad
\ y \ =\ \tfrac43x+\tfrac23\,.
$$
Indeed, we could have immediately seen that the slope is the vertical velocity 
over the horizontal velocity: $m=\frac43$.
\\[1em]
{\bf Parametric circles.} Given the unit circle defined by the equation $x^2+y^2 =1$,
we would like to {\it parametrize} it: to trace the curve by a particle moving
according to $(x(t),y(t))$.
One way is to let the particle make an angle of $t$ radians at time $t$, meaning:
$$
(x(t), y(t)) = (\cos(t), \sin(t)) \quad\text{for}\quad t\in[0,2\pi]\,.
$$
\\[-.5em]
\centerline{
\includegraphics[width=2.6in]
{10-1-Circle.png} 
}
\vspace{-1.0em}

\noindent
If we keep the same motion for all $t$, the particle 
travels around and around the circle.
We can check that this formula does trace the circle, because
the coordinates do satisfy the known equation:
$$
x(t)^2 + y(t)^2 \ =\ \cos^2(t) + \sin^2(t) = 1\,.
$$
Our standard circular motion has center $(0,0)$, radius $r=1$, 
starting at $(1,0)$ for $t=0$, 
with 1 counterclockwise rotation during $t\in[0,2\pi]$.
We can modify each part of this:
\begin{itemize}
\item Stretch the radius to $r=5$:
\ 
$(x(t), y(t)) = (5\cos(t), 5\sin(t))$. 

\item Move the center of the circle to $(6,7)$: 
\ 
$(x(t), y(t)) = (6+\cos(t), 7+\sin(t))$. 

\item Start from the bottom point $(0,-1)$ at $t=0$:
$(x(t), y(t)) = (\cos(t{-}\frac\pi2), \sin(t{-}\frac\pi2))$. 

\item Make the rotation clockwise:\footnote{
In general, to reverse the motion of $(x(t),y(t))$, take $(x(-t), y(-t))$
making time go backwards.
}
\ 
$(x(t), y(t)) = (\cos(-t), \sin(-t))$. 

\item Do 10 rotations over $t\in[0,1]$: \
$(x(t),y(t))= (\cos(10{\cdot}2\pi t),\sin(10{\cdot}2\pi  t))$.
\\[-.5em]

\end{itemize}

\noindent  Combining all the above modifications (from the top of the list down):
$$
(x(t), y(t)) \ =\ \left(\,6+5\cos\!\left(-10 {\cdot}2\pi t {-} \tfrac\pi2\right)\!, 7+5\sin\!\left(-10 {\cdot}2\pi t {-} \tfrac\pi2 \right) \right)\,.
$$
{\sc example:} We parametrize an ellipse, which is a circle stretched horizontally and/or vertically. For example, here is a parametric equation for
the ellipse centered at $(0,0)$, with extreme points $(\pm 3,0)$ and $(0,\pm2)$:
$$
(x(t),y(t)) \ =\ (3\cos(t), 2\sin(t))
$$
Because of the uneven stretching, the particle will {\it not} travel at constant speed,
and the central angle at time $t$ will {\it not} be proportional to $t$.
\\[.5em]
{\sc example:} A tricky way to parametrize the unit circle
is the {\it rational parametrization} which we derived in \S7.3 
to reduce trig integrals to rational functions:
$$
(x(t), \, y(t)) = \left(\frac{1-t^2}{1+t^2}\ , \ \frac{2t}{1+t^2}\right)\,.
$$
We can tell that each point $(x(t),y(t))$ lies on the circle because:
$$\begin{array}{rcl}
x(t)^2 + y(t)^2  &=&\ds
\left(\frac{1-t^2}{1+t^2}\right)^{\!2}
+ 
\left(\frac{2t}{1+t^2}\right)^{\!2}
\\[1.5em]
&=&\ds
\frac{1 - 2t^2 + t^4}{(1+t^2)^2}
 + \frac{4t^2}{(1+t^2)^2}
\\[1.5em]
&=&\ds
\frac{1 + 2t^2 + t^4}{(1+t^2)^2}
\ =\ 
\frac{(1+t^2)^2}{(1+t^2)^2}
\ =\ 1\,.
\end{array}$$
It is not easy to find any points on the circle with both coordinates
rational numbers: if we start with rational $x$, 
we expect irrational \smash{$y=\sqrt{1{-}x^2}$}. 
Amazingly, the rational parametrization 
produces infinitely many rational points: 
just plug in any fraction for $t$, for example 
$(x(\tfrac12),\,y(\tfrac12)) = (\tfrac35, \tfrac45)$
and
$(x(\tfrac23),\,y(\tfrac23)) = (\tfrac{5}{13}, \tfrac{12}{13})$.
For each point $(x,y)=(\frac ac, \frac bc)$,
clearing denominators in the equation $x^2+y^2=1$
gives $a^2+b^2=c^2$. This defines a {\it Pythagorean triple}:
a right triangle in which all three sides $a,b,c$ are whole numbers! 
For example $a=5$, $b=12$, $c=13$ satisfy $5^2+12^2 = 13^2$.
This is a sample of the mathematical field of Algebraic Geometry.
\\[1em]
{\bf Cycloid curve.} This famous curve 
 traces the path of a particle on the rim of a rolling wheel
(a unit circle rolling over the $x$-axis at unit speed).
\\[1em]
\centerline{
\includegraphics[width=4.5in]
{10-1-Cycloid.png} 
}
\vspace{0em}

\noindent
As the wheel rolls by $t$ radians, its circumference traces an 
equal distance $t$ along the $x$-axis, 
so the wheel's center moves to $(t,1)$.
If the center were fixed at the origin, the particle
on the rim would start at $(-1,0)$ and
turn clockwise once over $t\in[0,2\pi]$,
so its position would be:
$$
(\cos(-t{-}\tfrac\pi2), \sin(-t{-}\tfrac\pi2))
\ =\
(-\sin(t), -\cos(t))\,.
$$ 
Adding the linear motion of the center and the circular motion
around the center gives the parametric equation of the cycloid curve:
$$
(x(t), y(t)) \ =\
(t-\sin(t), \, 1-\cos(t))\,.
$$
These equations allow us (or a computer) to easily plot the cycloid.

Let us deparametrize this to get an $xy$-equation for the cycloid.
We solve for $t$ in terms of one variable (in this case $y$), and plug into the other variable (in this case $x$):\footnote{Strictly, $\sin(t)=\pm\sqrt{1-\cos^2(t)}$, but the minus sign just gives an extraneous upside-down cycloid.} 
$$\left\{ \begin{array}{c}
x = t-\sin(t)
\\
y = 1-\cos(t)
\end{array}\right.
\quad\Longrightarrow\quad
\left\{ \begin{array}{rcl}
\cos(t) &=& 1-y
\\[.3em]
\sin(t) &=& \sqrt{1{-}(1{-}y)^2} \ = \ \sqrt{2y{-}y^2}
\\[.3em]
t &=& \arccos(1{-}y)
\\[.3em]
x &=& \arccos(1{-}y)-\sqrt{2y{-}y^2}
\end{array}\right.
$$
%There is some ambiguity with $\arccos$ because $\cos$ is not one-to-one. In fact, given the standard values $0\leq \arccos(z)\leq \pi$,
%the full cycloid is given by a union of graphs $x=f(y)$ for all integers $n$:
%$$
%x \ = \ 2\pi n \pm \left(\arccos(1{-}y)-\sqrt{2y{-}y^2}\,\right)
%\quad\text{for $y\in[0,2]$,}
%$$
Simplifying:
$$
\cos\!\left(x+\sqrt{y(2{-}y)}\right) + y \ =\ 1\,.
$$
This weird $xy$-equation lets us easily check if {\it a given point} $(x,y)$ lies on the cycloid.
The parametric form, on the other hand, allows us to {\it produce points} on the curve.
\\[1em]
{\bf Epicycloids.} 
One variant of the cycloid is the {\it epicycloid}, in which
the wheel rolls around a fixed circle. The curve varies 
depending on the relative size of the two circles. 
From the perspective of a fixed central Earth, the trajectories 
of the other planets are very close to epicycloids,
and the classical astronomers in the tradition of Ptolemy 
attempted to find an exact model for planetary motion by adding 
further epicycles, wheels rolling on wheels like a gigantic clockwork.
But starting with Copernicus, we interpret the apparent
epicycloid as an illusion based on combining
the separate orbits of Earth and the other planet around a 
fixed central Sun.

Here is a compound epicycloid with a central circle of radius 1,
a wheel of radius $\frac16$ rolling around it, and a wheel of radius
$\frac16$ rolling around that (assuming the circles can pass through each other):%\ $(x(t), y(t)) = (\frac76\cos(\frac16 t) + \frac12\cos(2t), \ \frac76\sin(\frac16 t) + \frac12\sin(2t))$.
\\[1em]
\centerline{
\includegraphics[width=4.5in]
{10-1-Epicycloid.png} 
}
%\\[1em]
%{\bf Cusp curve.} One more example: $(x(t),y(t))=(t^3,t^2)$. This most easily pictured if we deparametrize: $t=x^{1/3}$ and $y=(x^{1/3})^2 = x^{2/3}$,
%which looks like this:
%\\[1em]
%\centerline{
%\includegraphics[width=3in]
%{10-1-Cusp.png} 
%}
%\vspace{0em}
%
%\noindent
%Here the tangent slopes are given by: $\frac{dy}{dx} = \frac23 x^{-1/3}
%= \frac{2}{3\sqrt[3]x}$, and at $x=0$ this becomes $\pm\infty$, a vertical tangent.
%Looking at this parametrically, the tangent line at $t=0$ is:
%$(x(0)+x'(0)t, y(0) + y'(0)t) = (0+0t, 0+0t)$. That is, the particle stops momentarily
%at the cusp point before starting in the opposite direction: it bounces.



\end{document}
%%%%%%%%%%%  Picture  %%%%%%%%%%
\\[1em]
\centerline{
%\includegraphics[width=2in]
{1-8-Discont-iv.png}
\qquad \qquad 
%\includegraphics[width=2in]
{1-8-Discont-v.png} 
}
\vspace{0em}

\noindent
%%%%%%%%%%%%%%%%%%%%%%%%%

%%%%%%%%%%%  Table  %%%%%%%%%%
In fact, we get the following derivatives:
$$\begin{array}{|c||c|c|c|c|c|c|}
\hline &&&&&& \\[-.9em] 
f(x) & \sin(x) & \cos(x) & \tan(x)  & \sec(x) & \csc(x) & \cot(x)
\\[.2em] \hline &&&&&& \\[-.9em]
f'(x) &\cos(x) & -\sin(x) & \sec^2(x) & \tan(x)\sec(x) & -\cot(x)\csc(x) & -\csc^2(x)
\\[.1em] \hline
\end{array}$$
\\[1em]
%%%%%%%%%%%%%%%%%%%%%%%%%%



















